%% institucion.tex — Datos de la universidad y formato de las portadas.
%% Un comando por línea. Todos son opcionales: si uno queda vacío, { }, no
%% sale en las portadas. Normalmente no hace falta tocarlo: lo trae el perfil
%% de tu escuela (aquí, ESI de la UCLM).

\universidad{Universidad de Castilla-La Mancha}
\escuela{Escuela Superior de Informática}
\logo{esi_logo.pdf}                       % archivo dentro de estilo/ (.pdf, .png o .jpg); vacío = sin logo
\tipoTrabajo{tfg}                         % tfg | tfm | tesis | otro
\nombreTrabajo{}                          % (opcional) sustituye a «Trabajo Fin de Grado»…
\titulacion{Grado en Ingeniería Informática}
\etiquetaEspecialidad{Tecnología específica}  % vacío = no se muestra la especialidad (\intensificacion de datos.tex)
\idiomaPortadas{espanol}                  % espanol | ingles | documento (el de \idioma en datos.tex)

%% Formato de página
\tamanoLetra{12pt}                        % 10pt | 11pt | 12pt
\margenes{35mm}{20mm}{25mm}{25mm}         % interior (encuadernación), exterior, superior, inferior
\interlineado{1.5}                        % 1 | 1.15 | 1.25 | 1.5 | 2
